\section*{第 \ref{ch:b2:mass-spec-atomic} 章——质谱法与原子光谱法}

\begin{solution}{exo:b2:mass-spec-atomic:1}
78; $400/2 = 200$; $194 + 1 = 195$, $[\mathrm M + \mathrm H]^+$.
\end{solution}

\begin{solution}{exo:b2:mass-spec-atomic:2}
奇数质量：分子中氮原子数为奇数，对于如此小的分子只能含一个氮。\ce{C3H7NO}（丙酰胺，73）
或 \ce{C4H11N}（1-丁胺，73）。
\end{solution}

\begin{solution}{exo:b2:mass-spec-atomic:3}
近乎等高：一个溴。3 : 1：一个氯。
\end{solution}

\begin{solution}{exo:b2:mass-spec-atomic:4}
黄色：钠，\qty{589.0}{nm} 和 \qty{589.6}{nm}。深红色：锂，\qty{670.8}{nm}。
% ledger: asd:NaD2, asd:NaD1, asd:Li671
\end{solution}

\begin{solution}{exo:b2:mass-spec-atomic:5}
$n_C \approx 6.7/1.11 \approx 6$.
\end{solution}

\begin{solution}{exo:b2:mass-spec-atomic:6}
\ce{CO}：27.9949；\ce{N2}：28.0061；\ce{C2H4}：28.0313 u。分辨率：\ce{CO}/\ce{N2} $28/0.0112
\approx 2.5 \times 10^3$; \ce{N2}/\ce{C2H4} $28/0.0252 \approx 1.1 \times 10^3$; \ce{CO}/\ce{C2H4}
$28/0.0364 \approx 7.7 \times 10^2$.
% ledger: mass:C12, mass:O16, mass:N14, mass:H1
\end{solution}

\begin{solution}{exo:b2:mass-spec-atomic:7}
86：\ce{C5H10O} 的分子离子。57：\ce{C2H5CO+}，由 $\alpha$-裂解生成的酰基正离子（失去
一个乙基自由基，29）。29：\ce{C2H5+}。
% ledger: ms:pentan-3-one
\end{solution}

\begin{solution}{exo:b2:mass-spec-atomic:8}
86：$\mathrm M^{+\bullet}$。71：$\mathrm M - \ce{CH3}$，即酰基正离子 \ce{C3H7CO+}。43：\ce{CH3CO+}
（失去丙基自由基）。58：麦氏重排，经由丙基链的
$\gamma$ 碳上的一个氢，失去乙烯。在戊-3-酮中，乙基没有 $\gamma$ 碳：没有
麦氏重排离子；58 处的小峰是 57 离子的 \ce{^{13}C} 同位素峰（含三个碳，
约 3.3~\%）。
% ledger: ms:pentan-2-one, ms:pentan-3-one
\end{solution}

\begin{solution}{exo:b2:mass-spec-atomic:9}
过原点的直线：斜率 $\sum c_iA_i/\sum c_i^2 = 1.553/30.0 \approx \qty{0.0518}{L/mg}$。样品：
$0.128/0.0518 \approx \qty{2.47}{mg/L}$ 的铜。
\end{solution}

\begin{solution}{exo:b2:mass-spec-atomic:10}
$t = L\sqrt{m/(2eU)} = 1.00\sqrt{1000 \times 1.6605 \times 10^{-27}/(2 \times 1.6022 \times 10^{-19}
\times 2.00 \times 10^4)} \approx \qty{16.1}{\micro s}$。对于 1001 u，$t$ 的增大倍数为
$\sqrt{1.001} \approx 1 + 0.0005$: $\Delta t \approx \qty{8.0}{ns}$.
% ledger: const:u, const:e
\end{solution}

\begin{solution}{exo:b2:mass-spec-atomic:11}
两个氯：$(0.758 + 0.242x)^2$ 给出 $0.575$、$0.367$、$0.0586$：$m/z$ 为 84、86 和 88，比值为
100 : 64 : 10.
\end{solution}

\begin{solution}{exo:b2:mass-spec-atomic:12}
\ce{^{40}Ar^{16}O+}: $39.9624 + 15.9949 = 55.9573$; \ce{^{56}Fe}: 55.9349; $R = 56/0.0224 \approx 2.5 \times
10^3$。否则测量铁的另一同位素 \ce{^{57}Fe}，或者
在进入分析器之前用碰撞池去除 \ce{ArO+}。
% ledger: mass:Ar40, mass:Fe56, mass:O16
\end{solution}

\begin{solution}{pb:b2:mass-spec-atomic:1}
\textbf{1.} 156、158 和 160（以及 157 和 159）：该分子存在多种同位素形式。
\textbf{2.} 三个峰相隔两个质量单位，中间的峰最大：含一个氯和一个溴。
\textbf{3.} 156.
\textbf{4.} 偶数质量：不含氮，或含偶数个氮。
\textbf{5.} $156 - 79 - 35 = 42$: \ce{C3H6}.
\textbf{6.} \ce{C3H6BrCl}；不饱和度 $(2 \times 3 + 2 - 6 - 2)/2 = 0$：无环，无双
键。
\textbf{7.} $0.5/14.5 \approx 3.4~\%$，$3.4/1.11 \approx 3$ 个碳。
\textbf{8.} 157 处的峰高约为 0.5，只精确到 0.1：不确定度为 20~\%。
\textbf{9.} \ce{C3H6^{81}Br^{35}Cl} 和 \ce{C3H6^{79}Br^{37}Cl}。
\textbf{10.} $3 \times 12 + 6 \times 1.007825 + 78.918338 + 34.968853 \approx 155.9341$.
\textbf{11.} $156/(156.151 - 155.934) \approx 7.2 \times 10^2$.
\textbf{12.} 是：少两个氢，为 154。
\textbf{13.} 失去一个溴原子（79 或 81）：\ce{C3H6Cl+}，77（含 \ce{^{35}Cl}）、79（含 \ce{^{37}Cl}），
3 : 1.
\textbf{14.} 失去 HBr：自由基阳离子 $\mathrm{C_3H_5Cl^{+\bullet}}$，偶数质量。
\textbf{15.} \ce{C3H5+}，即烯丙基正离子，由两种卤素（先失 Br，后失 HCl）均失去后生成。
\textbf{16.} \ce{CH2Cl+}（49 和 51，3 : 1）。
\textbf{17.} \ce{CH2Br+}（93 和 95）和 \ce{C2H4Br+}（107 和 109），均为 1 : 1：溴位于链的一端，
第二个来自失去 \ce{CH2Cl}。
\textbf{18.} \ce{C-Br} 键比 \ce{C-Cl} 键弱，且 \ce{Br} 是更好的离去
基团。
\textbf{19.} 否：127--131 处无峰；\ce{CH2Cl+} 和 \ce{CH2Br+} 表明每个卤素都位于末端
\ce{CH2}.
\textbf{20.} \ce{BrCH2CH2CH2Cl}，1-溴-3-氯丙烷。
\textbf{21.} 156/158/160: M; 77/79: M $-$ Br; 76: M $-$ HBr; 107/109: M $-$ \ce{CH2Cl};
93/95: \ce{CH2Br+}; 49/51: \ce{CH2Cl+}; 41: \ce{C3H5+}; 39, 27: \ce{C3H3+}, \ce{C2H3+}.
\textbf{22.} 否：不存在羰基。
\textbf{23.} 1-溴-2-氯丙烷，\ce{CH3CHClCH2Br}：失去 \ce{CH2Br} 得到 63/65 处的 \ce{CH3CHCl+}。
\textbf{24.} \ce{C3H6^{81}Br^{37}Cl}.
\textbf{25.} $M$: $0.758 \times 0.5069 = 0.384$; $M+2$: $0.758 \times 0.4931 + 0.242 \times 0.5069 =
0.496$；$M+4$：$0.242 \times 0.4931 = 0.119$。比值 $\boldsymbol{\approx 3.2 : 4.2 : 1}$，约为 3 : 4 :
1；谱图给出 $14.5 : 18.6 : 4.4 = 3.3 : 4.2 : 1$。
\end{solution}
